\section{Conclusion (Provisional, II-A Only)}
\label{sec:conclusion}

This report establishes, under a behaviorally validated offline replay
model of a real 30B-A3B MoE coding-workload trace
(Section~\ref{sec:replay-validation}), that a substantial gap exists
between causal LRU-family expert retention and a non-causal Belady-optimal
upper bound, at every tested capacity-constrained expert-memory budget
(Sections~\ref{sec:oracle-headroom}--\ref{sec:sensitivity-capacity}). We
decompose this gap into a small, non-monotonically-signed structural
component attributable to the runtime's sharded retention architecture,
and a substantially larger policy-information headroom
(Section~\ref{sec:recoverable-envelope}) -- the scientifically clean
upper bound motivating future causal retention and prediction research.
The largest tested relative headroom occurs near $8$~GiB ($70.66\%$) and
the largest tested absolute headroom near $1.5$~GiB
($\approx 103.57$~GB); the gap collapses entirely once the working set
becomes effectively resident at $16$~GiB, the first tested such point.

This is a provisional conclusion, limited strictly to Research II-A. It
establishes an upper bound on recoverable retention performance; it does
not establish, estimate, or imply any demonstrated causal recoverability,
latency improvement, or production performance gain. Research II-B
(whether exploitable routing structure exists) and Research II-C
(whether a causal predictor can recover any of this headroom, and
validation against the real runtime) remain future work, not yet
started. The present II-A tranche is frozen as of the commit and run IDs
recorded in Appendix~\ref{sec:provenance-appendix} and
\texttt{II-A-FREEZE.json}, and is intended to be incorporated, without
silent strengthening of its claims, into the eventual combined
II-A+II-B+II-C Research Part 2 paper.
